\defmodule {latl}

This program applies the spectral and/or lattice test(s) to a given
generator.  It uses the {\tt LONGINT} representation for the multipliers
and modulus (module {\tt MULTLI}) and the {\tt LONGREAL} representation
for the bases (module {\tt BASISLR}).
The generator analyzed can be a combined generator with $J$ components,
and expressed that way in the data file.
\iffalse  %%%%%%
If $J=1$ and $m$ is a power of two expressed as
$(x\ y\ z) = (2\ y\ 0)$ in the data file,
then the lattice associated to the multiplicative LCG with power
of two modulus, as described in \cite{FIS90}, is analyzed.
Otherwise, the lattice analyzed is the smallest one that contains all
the points, for all possible starting values (i.e. the union of all subcycles,
when there are subcycles).
\fi  %%%%%%
To run the program, type ``{\tt latl \<\em file\>}'', where \<{\em file}\> is
the name of the data file, without extension.

The data must be in a file with extension
``{\tt .dat}'', according to the format of Figure~\ref{lat.dat}.
The data fields have the same meaning as for {\tt seekl},
except that several fields have been removed, 
$a_i$ replaces $b_i$ and $c_i$,
there is a single category for the dimensions, 
and the options for {\em LatticeInfo} and {\em ResultForm\/} are different.
\begin {description}
\item [$(a_1,\ldots,a_k)$ :] is the vector of (integer) multipliers.
  They can be given in format (a) or (b) described in {\tt seekl}.
\item [{\em LatticeInfo} ~\<{\em Norm\/}\> :]
  If the value of {\em LatticeInfo} is {\tt Spectral} 
  or {\tt BeyerSpectral} or {\tt SpectralP},
  then the field {\em Norm\/} must appear.
  It indicates which type of normalization is used
  in the definition of $S_t$.
  The allowed values are {\tt (BestLat, Laminated, Rogers, Minkowski)}.
\item [{\em ResultForm} :]
  Selects in which form the results will be given.
  The possible values are {\tt (Terminal, RES, TEX)}.
  Lowercases are also allowed.
  \begin {description}
   \item [{\tt Terminal}] indicates that the results will be given on the
     terminal screen.
   \item [{\tt RES}] says that the results will be in a file with the same
     name as the data file, but with extension ``{\tt .res}''.
   \item [{\tt TEX}] asks the program to produce a file intended for
     \LaTeX, with extension ``{\tt .tex}''.
  \end {description}
\end {description}

%%  DATA FILE  %%%%%%%%%%%%%%%%%%%%%%%%%%%%%%%%%%%%%%%%
\setbox0=\vbox {\hsize=12.0cm \it \begin {tabular}{l}
           ReadGenFile ~$[$\<file0\>$]$ \\
           $J$ \\
           $\left.\begin {array}{l}
           \mbox {TypeGen} \\
           m \\
           k \\
           a_1 \\
           \vdots \\
           a_k \\
           \end {array} \ \ \ \right\} \mbox { (Repeat $J$ times) }$ \\
           Dim$(0)$ ~Dim$(1)$ \\
           LatticeInfo  ~\<{\em Norm\/}\> \\
           LatticeType \\
           LaGroupSizes ~Spacing\\
           VerifyBB \\
           MaxNodesBB \\
           ResultForm \\
\end {tabular}
}
\begin{figure}[htb] \centerline {\boxit{\box0}}
\caption { Data file format for latl. }  \label {lat.dat}
\end{figure}
